\chapter{Diseño de la solución}
\label{chap:diseno}


\section{Modelo conceptual del dominio}
\label{sec:modelo_dominio}

El modelo conceptual utilizado para el desarrollo de la solución representa la estructura estándar mediante la cual se implementa un componente compatible con el ecosistema Flamapy. En este caso, dicho modelo ha sido utilizado para representar un generador automático de modelos de características expresados en UVL, definiendo las entidades necesarias para configurar y ejecutar el proceso de generación. La solución se ha desarrollado siguiendo el paradigma de programación orientada a objetos de Python, donde cada entidad del modelo se implementa mediante clases organizadas en archivos con extensión \texttt{.py}. Esta estructura permite separar las responsabilidades de los diferentes componentes de la herramienta, diferenciando la configuración del generador de la operación encargada de producir los modelos de características.

La integración con UVLHub comienza con la recogida y comprobación de los
parámetros en los distintos pasos del asistente. Durante este recorrido,
UVLHub conserva la configuración como un diccionario en la sesión de Flask.
Al solicitar la generación, el navegador obtiene dicho diccionario en formato
JSON. Si no se solicita la comprobación de satisfacibilidad,
\texttt{fmgen\_wrapper.py} construye \texttt{FmgeneratorModel} y ejecuta
\texttt{GenerateFeatureModel} dentro de Pyodide. Si se activa dicha opción,
el navegador envía los parámetros al servidor, donde se generan y comprueban
los modelos booleanos mediante SAT. Los resultados se ofrecen para su
visualización o descarga, sin almacenarlos de forma permanente.

\subsection{Estructura del modelo}
\label{subsec:estructura_modelo}

El modelo conceptual representa la estructura interna del componente desarrollado siguiendo una arquitectura compatible con Flamapy. En este caso, dicho modelo ha sido adaptado para representar un generador automático de modelos de características expresados en UVL, separando la configuración del generador, la operación encargada de ejecutar la generación y el resultado producido. Véase la Figura~\ref{fig:modelado-conceptual}.

\begin{figure}[H]
    \centering
    \makebox[\textwidth][c]{%
        \includegraphics[width=1\textwidth]{figura_5_1_modelo_conceptual.pdf}
    }
    \caption{Modelado conceptual del sistema. Elaboración propia.}
    \label{fig:modelado-conceptual}
\end{figure}


Las entidades principales que componen el modelo son las siguientes:

\subsubsection*{\texttt{FmgeneratorModel}}

La entidad \textbf{FmgeneratorModel} representa la configuración global del generador de modelos de características. Actúa como entidad raíz del modelo conceptual y agrupa las diferentes configuraciones necesarias para definir las características del modelo que se desea generar, incluyendo aspectos relacionados con la nomenclatura, los niveles de expresividad UVL, la estructura jerárquica, las restricciones y los atributos. Esta entidad constituye la representación completa del estado del generador antes de iniciar el proceso de generación.

\subsubsection*{\texttt{NamingConfig}}

La entidad \textbf{NamingConfig} representa la configuración relacionada con la generación de nombres de los elementos creados durante el proceso de generación. Su finalidad es establecer las reglas utilizadas para definir la nomenclatura de los modelos y sus elementos asociados.

\subsubsection*{\texttt{LevelConfig}}

La entidad \textbf{LevelConfig} representa la configuración de los niveles de expresividad disponibles en el lenguaje UVL durante la generación. Esta entidad permite definir las capacidades del lenguaje que estarán habilitadas en los modelos generados, estableciendo las construcciones que podrán aparecer en el resultado final.

\subsubsection*{\texttt{FeaturesConfig}}

La entidad \textbf{FeaturesConfig} representa la configuración asociada a la generación de características del modelo. Su responsabilidad es definir las condiciones que controlan la creación de los elementos principales que forman parte de la estructura del modelo de características.

\subsubsection*{\texttt{HierarchyConfig}}

La entidad \textbf{HierarchyConfig} representa la configuración relacionada con la estructura jerárquica del modelo de características generado. Esta entidad define los parámetros necesarios para construir el árbol de características y establecer las relaciones jerárquicas entre sus elementos.

\subsubsection*{\texttt{ConstraintsConfig}}

La entidad \textbf{ConstraintsConfig} representa la configuración utilizada para la generación de restricciones entre los elementos del modelo. Su objetivo es definir cómo se crean las relaciones adicionales entre características y atributos, permitiendo generar modelos con diferentes niveles de complejidad y expresividad.

\subsubsection*{\texttt{AttributesConfig}}

La entidad \textbf{AttributesConfig} representa la configuración asociada a la generación de atributos vinculados a las características del modelo. Esta entidad permite establecer cómo se incorporan atributos al modelo generado, tanto mediante generación automática como mediante configuraciones definidas manualmente.

\subsubsection*{\texttt{GenerateFeatureModel}}

La clase \texttt{GenerateFeatureModel} representa la operación encargada
de generar un objeto \texttt{FeatureModel} a partir de la configuración
definida en \texttt{FmgeneratorModel}. Hereda de \texttt{Operation} y
dispone de un constructor sin parámetros. Su método \texttt{execute}
recibe la configuración, el índice del modelo y el número de intento,
valida los parámetros y coordina la generación de la jerarquía, los
atributos y las restricciones.

La configuración recibida se almacena en el atributo \texttt{model},
mientras que el objeto generado se guarda en \texttt{result}. El método
\texttt{execute} devuelve la propia operación, y el resultado se obtiene
mediante \texttt{get\_result}. Cada ejecución genera un único modelo de
características. La conversión a texto UVL y la asignación del nombre
del fichero se realizan fuera de esta clase.

\subsubsection*{\texttt{FeatureModel} generado}

El objeto \texttt{FeatureModel} representa el resultado de una ejecución
del generador. Contiene la raíz del modelo, la jerarquía de
características, sus relaciones, los atributos y las restricciones.
Constituye una representación en memoria compatible con el metamodelo
de características de Flamapy. Su transformación a texto UVL se realiza
posteriormente mediante \texttt{UVLWriter}, fuera de la operación
\texttt{GenerateFeatureModel}.

\subsection{Restricciones del dominio}
\label{subsec:constraints_dominio}



\begin{itemize}
    \item \textbf{Dependencia entre niveles de expresividad UVL:} El nivel
    booleano constituye la base de los modelos generados. En la configuración
    implementada, la activación del nivel de tipos requiere que también esté
    activo el nivel aritmético. Las opciones de los niveles secundarios se
    habilitan de acuerdo con los niveles principales seleccionados.

    \item \textbf{Validez sintáctica y satisfacibilidad:} La construcción de
    un objeto \texttt{FeatureModel} no demuestra por sí sola que su
    representación textual pueda ser interpretada por un lector UVL. Tampoco
    equivale a demostrar que el modelo sea satisfacible. La opción basada en SAT
    permite comprobar la satisfacibilidad de los modelos booleanos; la
    generación de modelos con niveles aritmético y de tipos proporciona
    artefactos con los que evaluar y desarrollar razonadores para esas
    construcciones. Cuando se activa la opción de satisfacibilidad, el sistema
    utiliza una configuración efectiva limitada a construcciones booleanas
    básicas; la comprobación SAT y los posibles reintentos se realizan en el
    servidor.

    \item \textbf{Reproducibilidad mediante el uso de semillas:} El usuario puede establecer la semilla utilizada para controlar la aleatoriedad del proceso de generación. Esta decisión permite repetir una ejecución utilizando la misma semilla y los mismos parámetros de configuración, siempre que se mantengan la versión del generador y las condiciones de ejecución. De esta forma, se facilita la reproducción de experimentos, la comparación de resultados y el análisis de posibles errores. Para ello, la clase \texttt{FmgeneratorModel} incorpora el atributo \texttt{seed} como parte de la configuración del generador.
\end{itemize}

\subsection{Espacio de soluciones y configuraciones}
\label{subsec:espacio_soluciones}


El espacio de configuraciones del generador está determinado por los valores que pueden tomar los diferentes parámetros configurables de la clase \texttt{FmgeneratorModel}. Para calcular el número de configuraciones posibles sería necesario revisar los límites entre los cuales puede variar cada parámetro. Debido a que algunos de ellos no presentan un límite superior establecido, como ocurre con el número de características, además de que existen parámetros numéricos reales, el espacio de configuraciones posibles es ilimitado.

Asimismo, el uso de mecanismos pseudoaleatorios permite explorar distintas soluciones para una misma configuración mediante la variación de la semilla. Esto favorece la diversidad de los modelos generados, aunque el uso de semillas diferentes no garantiza que todos los resultados sean distintos.

\section{Arquitectura del sistema}
\label{sec:arquitectura}

\subsection{Visión general}
\label{subsec:arquitectura_general}



La arquitectura diseñada organiza los diferentes componentes del sistema en capas lógicas, separando las responsabilidades asociadas a la interacción con el usuario, la gestión de la configuración, la generación de modelos y la infraestructura necesaria para su ejecución.

La \textbf{capa de presentación} está formada por la interfaz web proporcionada por UVLHub, mediante la cual el usuario configura los diferentes parámetros del generador. Esta interfaz sigue un diseño basado en un flujo tipo asistente, permitiendo organizar la configuración del modelo en diferentes pasos y facilitar la comprensión del proceso por parte del usuario.

La \textbf{capa de API/controladores} está formada principalmente por las \textit{routes} del módulo \textit{generator} de UVLHub. Estas funciones se encargan de gestionar la comunicación entre los diferentes pasos del flujo de configuración, enviando y recibiendo los parámetros introducidos por el usuario hasta completar la configuración necesaria para ejecutar el proceso de generación.

La \textbf{capa de dominio/lógica} está formada por la lógica interna del generador desarrollado siguiendo la arquitectura de extensiones de Flamapy, junto con los servicios encargados de gestionar la configuración dentro de UVLHub. El componente compatible con Flamapy contiene el modelo de configuración y la operación que produce los modelos de características. Las rutas y la lógica del asistente de UVLHub comprueban los parámetros y mantienen el diccionario de configuración en la sesión. En la generación ordinaria, \texttt{FmgeneratorModel} se construye y la operación se ejecuta en el navegador mediante Pyodide. En la generación satisfacible, el servicio del asistente construye el mismo modelo de configuración y ejecuta la operación en el servidor, donde se aplica la comprobación mediante PySAT.

La \textbf{capa de infraestructura} está formada por los elementos necesarios para la ejecución y despliegue de la aplicación. Para el desarrollo local se ha utilizado Docker junto con sus correspondientes contenedores y configuraciones, permitiendo construir y levantar el entorno completo de la aplicación de forma reproducible. Además, se ha utilizado Render como plataforma de despliegue principalmente para facilitar a los tutores la revisión del funcionamiento de la aplicación durante el desarrollo del proyecto. Asimismo, esta capa incluye el entorno de ejecución basado en WebAssembly mediante Pyodide, utilizado para ejecutar el generador en el navegador en el flujo ordinario, y las dependencias de PySAT utilizadas por el servidor cuando se solicita la comprobación de satisfacibilidad.

La funcionalidad de generación no añade persistencia para las configuraciones
ni para los modelos producidos. El diccionario de la sesión se utiliza de
forma temporal y los resultados se ofrecen al usuario para su visualización
o descarga. Esta descripción se refiere al generador integrado, sin excluir
los mecanismos de persistencia que UVLHub utiliza para otras funcionalidades.

\subsection{Diagrama de componentes}
\label{subsec:diagrama_componentes}


El diagrama de componentes mostrado en la Figura~\ref{fig:diagrama-componentes} representa la distribución de los principales módulos del sistema y la comunicación existente entre ellos.

El diagrama ofrece una vista estructural general; los recorridos de
generación ordinaria y satisfacible se detallan en las Figuras
\ref{fig:secuencia-generacion}, \ref{fig:secuencia-generacion-sat}
y \ref{fig:flujo_datos_principal}.

\begin{figure}[H]
    \centering
    \includegraphics[width=\textwidth]{diagrama_componentes.png}
    \caption{Diagrama de componentes del sistema. Elaboración propia.}
    \label{fig:diagrama-componentes}
\end{figure}

\subsection{Descripción secuencial del flujo de generación}
\label{subsec:descripcion_secuencial}


\begin{enumerate}

    \item \textbf{Inicio de la configuración:} El usuario accede al generador
    mediante la interfaz basada en asistente proporcionada por UVLHub y
    comienza a introducir los parámetros necesarios para definir las
    características del modelo que desea generar.

    \item \textbf{Comunicación entre los diferentes pasos del asistente:}
    Las \textit{routes} del módulo \textit{generator} reciben los parámetros
    introducidos por el usuario en cada paso del flujo de configuración y se
    encargan de enviarlos al siguiente paso junto con la información necesaria
    para continuar con la configuración del generador.

    En esta fase, la interfaz también realiza comprobaciones básicas, como la
    validación de las distribuciones de probabilidad, antes de enviar los
    valores al servidor.

    \item \textbf{Actualización y validación de los parámetros:}
    Tras cada paso, las rutas y los servicios internos de UVLHub comprueban los
    parámetros recibidos y actualizan el estado del asistente almacenado en la
    sesión de Flask. Estas comprobaciones permiten detectar valores inválidos
    antes de continuar. En caso de error, se informa al usuario para que
    corrija los parámetros correspondientes. Una vez completado el asistente,
    los valores se consolidan en un diccionario de parámetros; en esta fase
    todavía no se construye la instancia de \texttt{FmgeneratorModel}.

    \item \textbf{Obtención de la configuración y selección del flujo:}
    En el último paso del asistente, el navegador obtiene la configuración en
    formato JSON mediante el endpoint
    \texttt{/generator/random/params-json}. Si se solicita la generación
    ordinaria, prepara Pyodide y carga las dependencias necesarias para
    ejecutar el código Python mediante WebAssembly. Si se activa la opción de
    satisfacibilidad, envía los parámetros JSON mediante una petición
    \texttt{POST} a \texttt{/generator/random/generate-sat}.

    \item \textbf{Construcción de la configuración y ejecución del generador:}
    En el flujo ordinario, \texttt{fmgen\_wrapper.py} construye
    \texttt{FmgeneratorModel} mediante \texttt{from\_flat\_dict} dentro de
    Pyodide. Después ejecuta \texttt{GenerateFeatureModel} para cada modelo
    solicitado, obtiene el objeto \texttt{FeatureModel} y lo transforma en
    texto UVL mediante \texttt{UVLWriter}. En el flujo de satisfacibilidad,
    el servicio del asistente construye la configuración efectiva booleana
    en el servidor y genera cada candidato con la misma operación. Cada
    candidato se comprueba mediante PySAT; si no es satisfacible, se realiza
    otro intento, hasta alcanzar el límite establecido. Solo los modelos
    satisfacibles se serializan y devuelven al navegador. Los errores de
    preparación, comprobación o generación se comunican al usuario.

    \item \textbf{Visualización y descarga del resultado:} Los resultados
    generados correctamente se muestran en la interfaz de UVLHub para que el
    usuario pueda consultar su contenido UVL o descargarlos. La interfaz no
    presenta como válido ningún candidato cuya generación o comprobación
    haya fallado.

\end{enumerate}

La Figura~\ref{fig:secuencia-configuración} muestra el flujo de configuración.
Las Figuras~\ref{fig:secuencia-generacion} y
\ref{fig:secuencia-generacion-sat} muestran, respectivamente, la generación
ordinaria en el navegador y la generación satisfacible en el servidor.

\begin{figure}[!htbp]
    \centering
    \makebox[\textwidth][c]{%
        \includegraphics[width=1.3\textwidth]{diagrama_secuencia_generacion_1.png}
    }
    \caption{Diagrama de secuencia del flujo de configuración. Elaboración propia.}
    \label{fig:secuencia-configuración}
\end{figure}

\begin{figure}[!htbp]
    \centering
    \includegraphics[width=1.1\textwidth]{secuencia_generacion_ordinaria.png}
    \caption{Diagrama de secuencia de la generación ordinaria de modelos UVL
    en el navegador mediante Pyodide. Elaboración propia.}
    \label{fig:secuencia-generacion}
\end{figure}

\begin{figure}[!htbp]
    \centering
    \includegraphics[width=1.1\textwidth]{secuencia_generacion_sat.png}
    \caption{Diagrama de secuencia de la generación de modelos booleanos
    satisfacibles en el servidor mediante PySAT. Elaboración propia.}
    \label{fig:secuencia-generacion-sat}
\end{figure}

\FloatBarrier

\section{Patrones de diseño aplicados}
\label{sec:patrones}

\subsection*{Arquitectura por capas (\textit{Layered Architecture})}

La solución sigue una organización basada en capas lógicas, separando las
responsabilidades asociadas a la interfaz de usuario, comunicación con el
sistema, lógica de generación e infraestructura. Este patrón permite reducir
el acoplamiento entre componentes y facilitar la evolución del sistema,
permitiendo modificar una capa sin afectar al resto siempre que se mantengan
sus interfaces de comunicación.

\subsection*{Arquitectura basada en componentes independientes (Component-based Architecture)}

El generador desarrollado se estructura como un componente independiente compatible con el ecosistema Flamapy, utilizando los mecanismos y modelos proporcionados por el \textit{framework} para interactuar con él. Esta decisión permite separar la lógica de generación del resto de la aplicación, favoreciendo la reutilización, el mantenimiento y la evolución del sistema.

\subsection{Patrón de interacción estilo asistente (\textit{wizard})}
La interfaz de configuración del generador utiliza un patrón de interacción
basado en asistente, mediante el cual la introducción de parámetros se
divide en diferentes pasos independientes. Este patrón permite organizar una
gran cantidad de opciones de configuración de forma progresiva, facilitando
la comprensión del usuario y permitiendo realizar validaciones durante el
proceso de configuración. Además, se incorporó un panel de progreso que
permite consultar en todo momento el estado actual de la configuración
realizada.


\section{Diseño de la interfaz de usuario}
\label{sec:ui_ux}


La interfaz del asistente de configuración utilizada por el generador procede
del trabajo previo realizado por David Romero dentro de UVLHub. En particular,
la ampliación del asistente a seis pasos, el panel lateral de seguimiento de la
configuración y el selector visual de distribuciones fueron diseñados e
implementados con anterioridad a este TFG. Estas decisiones se describen en
esta sección porque forman parte de la interfaz desde la que se utiliza el
generador integrado, pero no constituyen una aportación propia de este trabajo.
Las contribuciones realizadas en el TFG se centran en la implementación del
componente generador compatible con Flamapy y en las adaptaciones necesarias
para integrarlo con UVLHub.

\subsection{Mapa de navegación}
\label{subsec:mapa_navegacion}


El mapa de navegación del generador sigue una estructura secuencial basada en
un flujo tipo asistente, donde la configuración del modelo de
características se divide en diferentes etapas independientes. Esta organización
permite separar los diferentes grupos de parámetros configurables según su
funcionalidad, facilitando al usuario la definición progresiva de la
configuración necesaria para realizar la generación del modelo.

El mapa de navegación y las capturas de los seis pasos se recogen en el
Apéndice~\ref{app:interfaz-generador}.

Las vistas principales del sistema son las siguientes:


\subsubsection*{\textit{Paso 1: Batch}}

Esta vista permite establecer los parámetros generales del proceso de generación,
incluyendo el número de modelos a generar, la semilla utilizada para controlar
la aleatoriedad y la nomenclatura de los modelos generados.


\subsubsection*{\textit{Paso 2: Language levels}}

Esta vista permite seleccionar los niveles de expresividad UVL que estarán
disponibles durante la generación, incluyendo la activación de los diferentes
niveles principales y secundarios soportados.


\subsubsection*{\textit{Paso 3: Feature tree}}

Esta vista permite configurar los parámetros relacionados con la generación de la
estructura del árbol de características, incluyendo el número de características,
la profundidad máxima y la distribución de relaciones jerárquicas.


\subsubsection*{\textit{Paso 4: Constraints}}

Esta vista permite definir los parámetros asociados a la generación de
restricciones entre características y atributos, incluyendo la distribución de
los diferentes tipos de restricciones y operadores disponibles.


\subsubsection*{\textit{Paso 5: Attributes}}

Esta vista permite configurar los parámetros relacionados con la generación de
atributos asociados a las características, incluyendo la posibilidad de generar
atributos aleatorios mediante una distribución de tipos configurable.


\subsubsection*{\textit{Paso 6: Output}}

Esta vista permite ejecutar el proceso de generación del modelo UVL utilizando
la configuración proporcionada en los pasos anteriores. Además, permite
activar la comprobación de satisfacibilidad, que aplica una configuración
efectiva limitada al nivel booleano, y seleccionar los sufijos utilizados en
los archivos generados.

Una vez finalizada la generación, el usuario puede visualizar el modelo obtenido
dentro de la plataforma o descargarlo para su posterior utilización.

La Figura~\ref{fig:fig:diagrama-navegacion-diseno} muestra una vista general acerca de la navegación a través del sistema.

\begin{figure}[!htbp]
    \centering
    \makebox[\textwidth][c]{%
        \includegraphics[width=1.2\textwidth]{diagrama_navegacion.png}
    }
    \caption{Mapa de navegación del generador UVL. Elaboración propia.}
    \label{fig:fig:diagrama-navegacion-diseno}
\end{figure}

\FloatBarrier

\subsection{Prototipo visual inicial}
\label{subsec:mockups}


Durante la fase inicial del desarrollo se realizó un prototipo de la interfaz
con el objetivo de cubrir los requisitos iniciales identificados para el
sistema. Este prototipo no representa completamente la versión final de la interfaz, ya que
durante el desarrollo se fueron incorporando diferentes mejoras y modificaciones
orientadas a mejorar la usabilidad y adaptar el sistema a las nuevas necesidades surgidas durante el desarrollo.

Sus cinco bocetos pueden consultarse en el Apéndice~\ref{app:prototipos-iniciales}.

\subsection{Decisiones de diseño visual}
\label{subsec:decisiones_visuales}


\begin{itemize}

    \item \textbf{Paleta de colores:} La interfaz mantiene la identidad visual proporcionada por
    UVLHub, utilizando principalmente tonos azules, blancos y grises para los diferentes
    elementos de la interfaz. Además, se utilizan colores específicos para representar
    diferentes estados del sistema, empleando rojo para mensajes de error, amarillo
    anaranjado para avisos y verde para operaciones completadas correctamente.

    \item \textbf{Tipografía:} Se mantiene la tipografía utilizada por UVLHub,
    basada en las familias \textit{Inter}, \textit{sans-serif} y \textit{Helvetica},
    utilizadas como estándares dentro de la plataforma, conservando la coherencia visual con el resto de la aplicación y proporcionando una
    lectura clara de los elementos de la interfaz.

    \item \textbf{Organización de la interfaz y flujo de configuración:} Debido a la
    cantidad de parámetros necesarios para configurar el generador, la interfaz utiliza
    un diseño basado en un flujo tipo asistente, separando la configuración en
    diferentes pasos independientes. Esta decisión permite reducir la complejidad
    percibida por el usuario y mejorar la comprensión del proceso de configuración.

    Además, dispone de un panel lateral de progreso que permite consultar en todo
    momento el estado actual del proceso, los pasos completados y un resumen de la
    configuración seleccionada sin necesidad de volver atrás durante el flujo.

    \item \textbf{Componentes visuales y sistema de diseño:} La interfaz utiliza los
    componentes proporcionados por el sistema de diseño utilizado en UVLHub, basado
    principalmente en Bootstrap. Además, la interfaz incorpora componentes específicos para el generador,
    como el selector de distribuciones probabilísticas utilizado para
    configurar la generación de modelos, permitiendo adaptar la interfaz a las
    necesidades particulares del sistema.

    \item \textbf{Ayudas visuales y gestión de la información:} Para facilitar la
    configuración de un sistema con una gran cantidad de parámetros, se muestran
    únicamente las opciones relacionadas con los niveles de expresividad seleccionados
    por el usuario. También se incluyen textos explicativos y \textit{tooltips}
    en determinados parámetros con el objetivo de proporcionar contexto adicional
    sobre su funcionamiento.

    \item \textbf{Responsive y adaptación visual:} Aunque la aplicación está orientada
    principalmente a su uso desde equipos de escritorio, la interfaz se ha diseñado
    utilizando una estructura adaptable que permite su correcta visualización en
    diferentes tamaños de pantalla.

\end{itemize}

\section{Flujo de datos principal}
\label{sec:flujo}


El flujo de datos del sistema comienza con los parámetros introducidos por el usuario mediante la interfaz basada en asistente. Las rutas del módulo \texttt{generator} de UVLHub reciben estos datos y, junto con los servicios internos, realizan las comprobaciones correspondientes a cada paso. La configuración se mantiene en forma de diccionario dentro de la sesión de Flask, permitiendo conservar los valores introducidos durante el recorrido por el asistente.

Una vez completada la configuración, el navegador obtiene estos parámetros en
formato JSON mediante el endpoint
\texttt{/generator/random/params-json}. A partir de este punto existen dos
recorridos. En la generación ordinaria, \texttt{fmgen\_wrapper.py} construye
\texttt{FmgeneratorModel} dentro de Pyodide mediante
\texttt{from\_flat\_dict} y ejecuta \texttt{GenerateFeatureModel}. Los objetos
\texttt{FeatureModel} obtenidos se transforman en texto UVL mediante
\texttt{UVLWriter} dentro del navegador. Cuando se activa la comprobación de
satisfacibilidad, el navegador envía los parámetros JSON mediante una petición
\texttt{POST} a \texttt{/generator/random/generate-sat}. El servidor construye
la configuración efectiva booleana, genera los candidatos, comprueba su
satisfacibilidad mediante PySAT y devuelve los modelos aceptados en formato
JSON. Ambos recorridos terminan en la interfaz, que permite visualizar y
descargar los resultados. Esta separación responde al requisito RNF-03.

La Figura~\ref{fig:flujo_datos_principal} representa el recorrido común de
los parámetros y las dos vías de generación.

\begin{figure}[!htbp]
    \centering
    \includegraphics[width=\textwidth]{flujo_datos_principal_actualizado.png}
    \caption{Flujo de datos de la generación ordinaria en Pyodide y de la
    generación satisfacible mediante SAT en el servidor. Elaboración propia.}
    \label{fig:flujo_datos_principal}
\end{figure}

\FloatBarrier

\subsection{Estructura de la representación intermedia}
\label{subsec:representacion_intermedia}


La representación intermedia de la configuración consiste en un diccionario de parámetros almacenado temporalmente en la sesión de Flask. Este diccionario permite mantener los valores introducidos por el usuario entre los diferentes pasos del asistente, sin construir todavía una instancia del modelo de dominio del generador.

Las rutas y los servicios del módulo \texttt{generator} gestionan y
comprueban los parámetros correspondientes a cada paso. El estado se conserva
por pasos en \texttt{session["wizard"]} y la configuración consolidada se
almacena en \texttt{session["params"]}. Cuando el navegador solicita la
configuración, el servidor proporciona este último diccionario en formato
JSON. En la generación ordinaria, \texttt{fmgen\_wrapper.py} lo convierte en
una instancia de \texttt{FmgeneratorModel} mediante
\texttt{from\_flat\_dict} dentro de Pyodide. En la generación satisfacible,
el navegador envía el mismo diccionario al servidor y allí se construye la
configuración efectiva utilizada por la operación.

El Código~\ref{lst:load-step-state} muestra cómo el módulo recupera de la
sesión los valores guardados para un paso y los combina con los valores
predeterminados.

\begin{lstlisting}[language=Python,
caption={Lectura del estado de un paso desde la sesión de Flask.},
label={lst:load-step-state}]
def load_step_state(step: int, defaults: dict):
    wizard = session.get("wizard", {})
    saved = wizard.get(str(step), {})

    out = defaults.copy()
    out.update(saved)

    return out
\end{lstlisting}

Esta representación se utiliza para mantener el estado temporal del asistente
y transferir los parámetros al entorno donde se realiza la generación. No
constituye un registro permanente de las configuraciones utilizadas ni de los
modelos obtenidos.

\subsection{Mapeo entre niveles}
\label{subsec:mapeo}


El mapeo entre el diccionario de parámetros y el modelo interno de
configuración se realiza mediante el método \texttt{from\_flat\_dict} de
\texttt{FmgeneratorModel}. En la generación ordinaria, el método se ejecuta
desde \texttt{fmgen\_wrapper.py} en el navegador; cuando se solicita la
comprobación de satisfacibilidad, se ejecuta desde el servicio del asistente
en el servidor. Esta transformación permite pasar de la representación
utilizada durante el asistente a la representación orientada a objetos que
necesita el generador.

Algunos ejemplos de este mapeo son:

\begin{itemize}
    \item \textit{NUM\_MODELS} $\rightarrow$ \textit{num\_models}: número de
    modelos de características que serán generados.

    \item \textit{TYPE\_LEVEL} $\rightarrow$ \textit{levels.type\_level}: activación de las construcciones asociadas al nivel de tipos del lenguaje UVL.

    \item \textit{MIN\_FEATURES} / \textit{MAX\_FEATURES} $\rightarrow$
    \textit{features.min\_features} / \textit{features.max\_features}: rango
    utilizado para determinar el número de características del modelo generado.
\end{itemize}

Una vez construida la instancia de \texttt{FmgeneratorModel}, la configuración
se valida y se utiliza para ejecutar \texttt{GenerateFeatureModel}. En el
flujo ordinario, la construcción y la ejecución tienen lugar en el navegador;
en el flujo satisfacible, se realizan en el servidor con una configuración
efectiva booleana y cada candidato se comprueba mediante PySAT. Los objetos
\texttt{FeatureModel} aceptados se transforman posteriormente a su
representación textual en UVL.


Con la arquitectura del sistema, los modelos conceptuales, los flujos de datos
y las principales decisiones de diseño definidos, queda establecida la solución
desde un punto de vista estructural y conceptual. El siguiente capítulo aborda
la implementación de dicha solución, detallando las tecnologías utilizadas y
cómo los diferentes componentes definidos en este diseño han sido desarrollados
para obtener el sistema final.